% Part: incompleteness
% Chapter: arithmetization-syntax
% Section: substitution

\documentclass[../../../include/open-logic-section]{subfiles}

\begin{document}

\olfileid{inc}{art}{sub}
\olsection{ప్రతిస్థాపన}

\tetoken{ఒక సూత్రం}{!!a{formula}}~$!A$లోని
\tetoken{ఒక చరం}{!!a{variable}}~$u$ యొక్క స్వేచ్ఛా ఘటనలన్నిటినీ పదం~$t$తో
మార్చే క్రియను ప్రతిస్థాపన అంటామని గుర్తుచేసుకోండి; దాన్ని
$\Subst{!A}{t}{u}$ అని రాస్తాం. \tetoken{చరాలు}{!!{variable}s},
\tetoken{సూత్రాలు}{!!{formula}s}, పదాల గ్యోడెల్ సంఖ్యలపై నిర్వహించినప్పుడు
ఈ క్రియ ఆదిమ పునరావృత్తం.

\begin{prop}\ollabel{prop:subst-primrec}
క్రింది ధర్మం గల ఆదిమ పునరావృత్త ప్రమేయం $\fn{Subst}(x,y,z)$ ఉంది:
\[
\fn{Subst}(\Gn{!A}, \Gn{t}, \Gn{u}) = \Gn{\Subst{!A}{t}{u}}.
\]
\end{prop}

\begin{proof}
ఆదిమ పునరావృత్తం ద్వారా $\fn{hSubst}$ ప్రమేయాన్ని ఇలా నిర్వచించవచ్చు:
\begin{multline*}
\begin{aligned}
\fn{hSubst}(x, y, z, 0) & = \emptyseq \\
\fn{hSubst}(x, y, z, i+1) & =
\end{aligned}\\
\begin{cases}
\fn{hSubst}(x, y, z, i) \concat y & \text{$\fn{FreeOcc}(x,z,i)$ అయితే} \\
\fn{append}(\fn{hSubst}(x, y, z, i), (x)_{i}) & \text{లేకపోతే.}
\end{cases}
\end{multline*}
ఇప్పుడు $\fn{Subst}(x,y,z)$ను $\fn{hSubst}(x,y,z,\len{x})$గా
నిర్వచించవచ్చు.
\end{proof}

\begin{prop}
\ollabel{prop:free-for}
గ్యోడెల్ సంఖ్య~$x$ గల సూత్రంలో, గ్యోడెల్ సంఖ్య~$z$ గల చరానికి గ్యోడెల్
సంఖ్య~$y$ గల పదం \tetoken{ప్రతిస్థాపనకు స్వేచ్ఛగా}{!!{free for}}
ఉంటేనే నిలిచే $\fn{FreeFor}(x,y,z)$ సంబంధం ఆదిమ పునరావృత్తం.
\end{prop}

\begin{proof} అభ్యాసం. \end{proof}

\begin{prob}
\olref[inc][art][sub]{prop:free-for}ను నిరూపించండి.
\end{prob}

\end{document}
